% !TeX root = ../../Book4.tex
% 纲要标题：## 附录 E　景与 topos 初步
% 纲要细目（写作范围）：
% > 用覆盖筛定义景与集合值层，建立预层、空间层和群作用的指数对象与子对象分类子；构造平展景并联系域上的 Galois 作用，一般景的层化和内射存在性给出证明概要。

\chapter{景与 topos 初步}
\label{b4:app:E}
\BookChapterNumber{b4:app:E}

\begin{chapterabstract}
覆盖可以由开集组成，也可以由代数对象间的态射组成。景将它们写成统一的局部拼接条件；集合值层范畴中的指数对象表示函数，子对象分类子则表示子对象。我们在预层、拓扑空间上的层和群作用中核验这些构造，并区分初等 topos、Grothendieck topos 与交换群值层范畴。平展覆盖把可分扩张也纳入局部研究；域上的平展层由连续 Galois 作用描述，截面的拼接由此转化为固定元素的计算。
\end{chapterabstract}

\begin{learninggoals}
\begin{itemize}
\item 由覆盖族写出其生成的筛，核验筛的拉回与拓扑公理。
\item 将层条件解释为相容限制的唯一拼接。
\item 构造预层的指数对象及筛分类子，核验它们的泛性质。
\item 比较集合值层与交换群值层，辨认两类 topos 的定义范围。
\item 核验平展覆盖的换基与复合，借有限平展对象计算域上的层与全局截面。
\end{itemize}
\end{learninggoals}

取一个开区间上的连续函数。它可能在较小区间上恒为零，在整个区间上却不是零。若用“是否为零”给这个函数附一个真假值，就会忽略限制以后发生的变化。把所有使它为零的开区间合起来，才能记录其局部性质。这正是子对象分类子在层范畴中所保存的信息：一个条件在哪些局部区域成立。

\ProseParagraphBreak
另外，在实数域上方程 \(T^2+1=0\) 没有解，在复数域上却有解。这次局部变化是 \(\operatorname{Spec}\mathbb C\to\operatorname{Spec}\mathbb R\)，而不是在单点空间里缩小开集。两份这样的局部变化怎样交叠，要由 \(\mathbb C\otimes_{\mathbb R}\mathbb C\) 描述。\secref{b4:sec:C04}已经从域扩张计算了相应的 Galois 作用；\secref{b4:sec:E04}将把这项计算放回景与层的定义中。

\ProseParagraphBreak
本附录使用本卷的泛性质、Yoneda 引理、有限极限与笛卡尔闭结构，平展覆盖还用到第三卷第23.3节的环平展判据和第24章的概形纤维积。一般景的层化构造给出概要，具体空间上的结论则直接证明。主要参考为\cite[Chapters I--IV]{MacLaneMoerdijk1992Sheaves}。

% !TeX root = ../../Book4.tex
% 纲要标题：### E.1 覆盖、景与集合值层
% 纲要细目（写作范围）：
% #### E.1.1 筛与 Grothendieck 拓扑
% #### E.1.2 层的相容族与唯一拼接
% #### E.1.3 开集覆盖与群作用的例子

\section{覆盖、景与集合值层}
\label{b4:sec:E01}

在拓扑空间上，局部数据的相容性由开集的交来检验。在域上，添入一个可分根也可以成为局部变换，却不是原空间中的开子集。要同时描述这两种情形，须把“哪些态射足以覆盖一个对象”作为数据。筛把覆盖连同它的全部限制放在一起，使定义不要求范畴中处处存在纤维积。

\subsection{筛与 Grothendieck 拓扑}
\label{b4:subsec:E01:01}

\begin{definition}\label{b4:def:sieve}
设 \(\mathcal C\) 为小范畴，\(U\in\mathcal C\)。一族以 \(U\) 为目标的态射 \(S\)，若满足：\(f:V\to U\) 属于 \(S\) 时，每个可复合的 \(g:W\to V\) 都有 \(f\circ g\in S\)，则称 \(S\) 为 \(U\) 上的\term[shai]{筛}[sieve]。等价地，它是可表预层 \(h_U=\operatorname{Hom}_{\mathcal C}(-,U)\) 的子预层。对 \(f:V\to U\)，定义
\[
 f^*S=\{\,g:W\to V\mid f\circ g\in S\,\}.
\]
这称为筛沿 \(f\) 的拉回。
\end{definition}
\symidx{\(f^*S\)：筛沿态射的拉回}

由一族箭头 \((u_i:U_i\to U)\) 生成的筛，包括所有能经过某个 \(u_i\) 分解的箭头。因而它记录的不只是这组箭头，还包括由它们继续限制得到的局部数据。最大筛为整个 \(h_U\)，其中包含 \(1_U\)。

\ProseParagraphBreak
例如，把链 \(0<1<2\) 看成偏序范畴。在对象 \(2\) 上，由 \(1\to2\) 生成的筛还必须包含 \(0\to2\)，因为后者经过 \(0\to1\to2\) 分解；它不包含 \(1_2\)，因而不是最大筛。沿 \(1\to2\) 拉回时，\(1_1\) 与 \(0\to1\) 都在拉回筛中，所得却是对象 \(1\) 上的最大筛。筛的封闭条件把一张箭头及它的全部进一步限制一起保留。

\begin{definition}\label{b4:def:grothendieck-site}
设 \(\mathcal C\) 为小范畴。对每个 \(U\in\mathcal C\)，指定一族筛 \(J(U)\)，并满足以下条件：
\begin{enumerate}
\item 最大筛属于 \(J(U)\)。
\item 若 \(S\in J(U)\)、\(f:V\to U\)，则 \(f^*S\in J(V)\)。
\item 若 \(S\in J(U)\)，且 \(R\) 为 \(U\) 上的筛，对每个 \(f:V\to U\) 属于 \(S\) 都有 \(f^*R\in J(V)\)，则 \(R\in J(U)\)。
\end{enumerate}
称 \(J\) 为 \(\mathcal C\) 上的\term[Grothendiecktuopu]{Grothendieck 拓扑}[Grothendieck topology]，\((\mathcal C,J)\) 为\term[jing-site]{景}[site]，\(J(U)\) 中的筛为覆盖筛。
\end{definition}
\symidx{\(J(U)\)：Grothendieck 拓扑指定的覆盖筛族}

第二条说覆盖在改变观察对象后仍是覆盖；第三条说先覆盖、再在每块上覆盖，能够组成覆盖。这种“拓扑”指定的是态射的覆盖关系，不是给对象集合添上一族普通开集。定义及一般构造见 Mac Lane 与 Moerdijk\cite[Chapter III, Sections 1--2]{MacLaneMoerdijk1992Sheaves}。

\ProseParagraphBreak
三条公理还推出覆盖筛对包含关系向上封闭：若 \(S\in J(U)\) 且 \(S\subseteq R\subseteq h_U\)，则 \(R\in J(U)\)。事实上，对每个 \(f:V\to U\) 属于 \(S\)，也有 \(f\in R\)，故 \(1_V\in f^*R\)。筛中一旦含恒等箭头，就包含以该对象为目标的全部箭头，所以 \(f^*R=h_V\) 是覆盖筛。第三条公理随即给出 \(R\in J(U)\)。增加已有覆盖的箭头及其限制，不会破坏覆盖性质。

\subsection{层的相容族与唯一拼接}
\label{b4:subsec:E01:02}

\begin{definition}\label{b4:def:site-set-sheaf}
设 \((\mathcal C,J)\) 为小景，\(F:\mathcal C^{\mathrm{op}}\to\mathbf{Set}\) 为预层。若对每个 \(U\) 及每个 \(S\in J(U)\)，限制映射
\[
 F(U)\longrightarrow\operatorname{Nat}(S,F),\qquad
 x\longmapsto\bigl(f\longmapsto F(f)(x)\bigr)
\]
为双射，则称 \(F\) 为该景上的\term[jihezhiceng]{集合值层}[sheaf of sets]。这些层及其自然变换组成的范畴记为 \(\operatorname{Sh}_{\mathbf{Set}}(\mathcal C,J)\)。
\end{definition}
\symidx{\(\operatorname{Sh}_{\mathbf{Set}}(\mathcal C,J)\)：景上的集合值层范畴}

右边的一项自然变换，就是对每个 \(f:V\to U\) 属于 \(S\) 指定 \(x_f\in F(V)\)，且有
\[
 x_{f\circ g}=F(g)(x_f).
\]
双射的单射性表示局部限制唯一决定整体，满射性表示相容的局部数据能够拼接。\(\mathcal C\) 小保证这里的自然变换构成集合。

\ProseParagraphBreak
若 \(S\) 由覆盖箭头 \(u_i:U_i\to U\) 生成，且各纤维积 \(U_i\times_UU_j\) 存在，那么筛上的相容族等价于一族 \(x_i\in F(U_i)\)，满足
\[
 F(p_1)(x_i)=F(p_2)(x_j),
\]
其中 \(p_1,p_2\) 是两张投影，等式在 \(F(U_i\times_UU_j)\) 中成立。自然变换沿 \(u_ip_1=u_jp_2\) 限制，便给出这个等式。反过来，对 \(f:V\to U\) 属于 \(S\)，选分解 \(f=u_i g\)，令 \(x_f=F(g)(x_i)\)。若还可写为 \(f=u_j h\)，纤维积的泛性质给出 \(V\to U_i\times_UU_j\)，把上式沿它限制就得到 \(F(g)(x_i)=F(h)(x_j)\)，所以定义与分解的选择无关；继续前复合也保持这一定义。这里的条件包括 \(i=j\)：同一覆盖箭头的不同分解也必须相容。

\begin{proposition}\label{b4:prop:open-site-sheaves}
设 \(X\) 为拓扑空间，以开集及包含为范畴 \(\mathcal O(X)\)。规定 \(U\) 上的筛覆盖，当且仅当其中所有开集的并为 \(U\)。这给出 Grothendieck 拓扑；集合值层条件等价于通常开覆盖上的相容截面唯一拼接条件。
\end{proposition}
\symidx{\(\mathcal O(X)\)：拓扑空间的开集范畴}
\begin{proof}
筛是在包含关系下向下封闭的开集族。最大筛显然覆盖。沿 \(V\subseteq U\) 拉回筛时，其中的开集覆盖 \(V\)，因为可用 \(V\cap W\) 限制原覆盖中的 \(W\)。若 \(S\) 覆盖 \(U\)，且 \(R\) 在每个 \(W\in S\) 上的拉回覆盖 \(W\)，则 \(R\) 中开集的并包含每个 \(W\)，故等于 \(U\)。这核验三条公理。

开集包含的纤维积为 \(U_i\cap U_j\)，所以前述相容等式就是截面在交集上的限制相同。任意覆盖筛本身也是开覆盖，故筛上的唯一拼接与开覆盖上的唯一拼接相互推出。空集由空族覆盖，还得到 \(F(\varnothing)\) 是单元素集合。
\end{proof}

第三卷第24.1--24.2节的局部化构造就是这一条件的具体实例。设 \(A\) 为交换环、\(M\) 为 \(A\) 模，\(X=\operatorname{Spec}A\)。结构层与伴随模层在基本开集上的截面分别为 \(\mathcal O_X(D(f))=A_f\) 与 \(\widetilde M(D(f))=M_f\)，限制由继续局部化给出。若 \(X=\bigcup_iD(f_i)\)，在 \(D(f_i)\cap D(f_j)=D(f_if_j)\) 上限制相同的局部分式便唯一黏合成全局截面。把这组基本开覆盖及其所有开子集放入生成的筛，就得到上一定义中的相容族；环乘法与模作用是这些集合值层上的附加结构。

\ProseParagraphBreak
这里取值于集合。\secref{b4:sec:C03}中的层取值于交换群，截面另有加法；不能仅因同样满足拼接就把两个层范畴当作同一个范畴。

\subsection{开集覆盖与群作用的例子}
\label{b4:subsec:E01:03}

若每个对象只有最大筛被规定为覆盖，则任意预层都是层，因为 Yoneda 引理给出 \(\operatorname{Nat}(h_U,F)\cong F(U)\)。这称为平凡的 Grothendieck 拓扑。若 \(G\) 为群，把它看成单对象范畴，则每个非空筛都含全部箭头：含 \(g\) 就含 \(gg^{-1}=1\)，继而含每个箭头。因此在平凡拓扑下，层范畴就是该范畴上的预层范畴，也就是右 \(G\)-集范畴；用 \(g\cdot x=x\cdot g^{-1}\) 可转为左 \(G\)-集。

\ProseParagraphBreak
这个单对象范畴只使用 \(G\) 的抽象群结构。若 \(G\) 还带拓扑，要让它作用于离散集合时连续，须另外要求每个元素的稳定子群都开。离散集合上的作用连续，当且仅当每个轨道映射 \(g\mapsto x\cdot g\) 连续，也就是每个点的原像开。其各个非空纤维正是稳定子群的陪集；其中 \(x\) 的原像就是稳定子群本身，而平移保持开集，故这些纤维全开当且仅当稳定子群开。例如，固定素数 \(p\)，\(\mathbb Z_p\) 在自身底集合上的右平移作用 \(x\cdot g=x+g\) 是一个抽象右 \(\mathbb Z_p\)-集。但每个元素的稳定子群都是 \(\{0\}\)，它不含任何 \(p^r\mathbb Z_p\)，因而不是开子群；把该底集合赋予离散拓扑以后，这个作用便不连续。后面的 Galois 层所对应的是离散连续作用，不能直接用全部抽象 \(G_K\)-集替代。

\ProseParagraphBreak
\subsecref{b4:subsec:C04:02}的有限平展对象采用联合满射的覆盖。对有限连续 \(G_K\)-集，这类覆盖允许把等变函数逐块拼接。开覆盖与平展覆盖虽然来源不同，层公理所做的事相同：只有与限制相容的数据才允许成为一个整体。后面研究的是这些集合值层范畴共同具有的范畴结构。

% !TeX root = ../../Book4.tex
% 纲要标题：### E.2 子对象分类子与初等 topos
% 纲要细目（写作范围）：
% #### E.2.1 特征映射与子对象分类子
% #### E.2.2 预层范畴的指数对象
% #### E.2.3 筛所给出的分类子

\section{子对象分类子与初等 topos}
\label{b4:sec:E02}

一个子集 \(A\subseteq X\) 可以用特征函数 \(X\to\{0,1\}\) 表达。对于预层，一个截面可能在限制后才落入子预层，单个真假值便不足以记录全部信息。分类子将这种“在哪些限制下成立”也存为一个对象。

\subsection{特征映射与子对象分类子}
\label{b4:subsec:E02:01}

\begin{definition}\label{b4:def:subobject-classifier}
设 \(\mathcal E\) 为具有有限极限的局部小范畴，终对象记为 \(1\)。给定对象 \(\Omega\) 及态射 \(\mathrm{true}:1\to\Omega\)。若对每个单态射 \(m:A\to X\)，存在唯一态射 \(\chi_m:X\to\Omega\)，使下图为拉回方块，则称 \((\Omega,\mathrm{true})\) 为\term[ziduixiang-fenleizi]{子对象分类子}[subobject classifier]：
\[
\begin{tikzcd}
A\arrow[r]\arrow[d,"m"']&1\arrow[d,"\mathrm{true}"]\\
X\arrow[r,"\chi_m"']&\Omega.
\end{tikzcd}
\]
\(\chi_m\) 称为该子对象的特征映射。
\end{definition}

集合中的分类子为 \(\Omega=\{0,1\}\)，\(\mathrm{true}(*)=1\)。拉回就是 \(\chi_m^{-1}(1)\)，因而唯一的特征映射在 \(A\) 上取 \(1\)，其余地方取 \(0\)。单态射 \(A\to X\) 与它的像包含代表同一子对象，不需要先把 \(A\) 字面写成子集。

\begin{definition}\label{b4:def:elementary-topos}
设 \(\mathcal E\) 为局部小范畴。若它具有有限极限、是\cref{b4:def:cartesian-closed}所述的笛卡尔闭范畴，并具有子对象分类子，则称它为\term[chudeng-topos]{初等 topos}[elementary topos]。
\end{definition}

指数对象 \(G^F\) 表示“从 \(F\) 到 \(G\) 的函数”；分类子则表示“\(X\) 的一个子对象”。两者都用可表性表达，但不是同一种对象。定义与基本例子见\cite[Chapter IV, Section 1]{MacLaneMoerdijk1992Sheaves}。

\subsection{预层范畴的指数对象}
\label{b4:subsec:E02:02}

\begin{proposition}\label{b4:prop:presheaf-exponential}
设 \(\mathcal C\) 为小范畴，\(\widehat{\mathcal C}=\operatorname{Fun}(\mathcal C^{\mathrm{op}},\mathbf{Set})\)。对预层 \(F,G\)，令
\[
 (G^F)(U)=\operatorname{Nat}(h_U\times F,G).
\]
沿 \(a:V\to U\) 的限制由前复合 \(h_a\times1_F\) 给出。这定义一个预层，且为 \(\widehat{\mathcal C}\) 中的指数对象。因此预层范畴笛卡尔闭。
\end{proposition}
\begin{proof}
\(h_a\) 保持恒等与复合，前复合也如此，故确实得到预层。对 \(\theta\in(G^F)(U)\)、\(x\in F(U)\)，定义
\[
 \operatorname{ev}_U(\theta,x)=\theta_U(1_U,x).
\]
\(\theta\) 的自然性给出
\(G(a)\theta_U(1_U,x)=\theta_V(a,F(a)x)\)，这正是求值映射与限制相容的等式。

若 \(b:H\times F\to G\) 为自然变换，对 \(z\in H(U)\) 定义
\[
 \widetilde b_U(z)_V(f,x)=b_V(H(f)z,x)
 \quad(f:V\to U,\ x\in F(V)).
\]
\(b\) 的自然性使右边成为 \(h_U\times F\to G\) 的自然变换，也使 \(\widetilde b:H\to G^F\) 自然。求值后在 \((z,x)\) 上恢复 \(b_U(z,x)\)。反过来，若先给 \(c:H\to G^F\)，令 \(b=\operatorname{ev}\circ(c\times1_F)\)，则 \(c\) 与 \(c_U(z)\) 的自然性分别给出
\[
 \widetilde b_U(z)_V(f,x)
 =c_V(H(f)z)_V(1_V,x)
 =c_U(z)_V(f,x).
\]
两个构造互逆，并与 \(H\) 的态射相容，所以给出指数泛性质。
\end{proof}

\(G^F\) 通常不是逐对象函数集 \(G(U)^{F(U)}\)。它必须保存全部限制的相容性；公式中的 \(h_U\) 正是把这些条件一起纳入的工具。该构造见\cite[Chapter I, Section 6]{MacLaneMoerdijk1992Sheaves}。

\subsection{筛所给出的分类子}
\label{b4:subsec:E02:03}

\begin{theorem}\label{b4:thm:presheaf-topos}
设 \(\mathcal C\) 为小范畴。令 \(\Omega(U)\) 为 \(U\) 上所有筛组成的集合，限制沿态射取筛的拉回，\(\mathrm{true}_U(*)\) 为最大筛。则 \(\Omega\) 是预层范畴 \(\widehat{\mathcal C}\) 的子对象分类子，故 \(\widehat{\mathcal C}\) 是初等 topos。
\end{theorem}
\begin{proof}
筛的拉回满足 \(1_U^*S=S\) 与 \((a\circ b)^*S=b^*(a^*S)\)，故 \(\Omega\) 是预层。预层的有限极限逐对象计算；尤其单态射可识别为逐对象的子集包含。给定子预层 \(A\subseteq F\)，定义
\[
 \chi_U(x)=\{\,f:V\to U\mid F(f)x\in A(V)\,\}.
\]
\(A\) 对限制封闭，所以这确实是筛。\(\chi_V(F(a)x)=a^*\chi_U(x)\) 逐项成立，故 \(\chi\) 自然。\(\chi_U(x)\) 为最大筛，当且仅当 \(1_U\) 属于该筛，也就是 \(x\in A(U)\)。逐对象拉回便恢复 \(A\)。

若 \(\psi:F\to\Omega\) 也有这一拉回性质，则对每个 \(f:V\to U\)，
\[
 f\in\psi_U(x)
 \ \Longleftrightarrow\ f^*\psi_U(x)\text{ 为最大筛}
 \ \Longleftrightarrow\ F(f)x\in A(V).
\]
第一步使用筛对前复合封闭：\(f\) 属于筛恰好表示其拉回含 \(1_V\)。第二步使用 \(\psi\) 的自然性及拉回性质。因此 \(\psi=\chi\)。结合\cref{b4:prop:presheaf-exponential}即得结论。
\end{proof}

对于两个对象和一支非恒等箭头 \(a:V\to U\) 的范畴，\(U\) 上有三个筛：空筛、只含 \(a\) 的筛、最大筛。中间一个筛表达“在 \(U\) 上尚未成立，但限制到 \(V\) 后成立”。这说明为什么预层的真假值不必只有两个。

% !TeX root = ../../Book4.tex
% 纲要标题：### E.3 空间、群作用与 Grothendieck topos
% 纲要细目（写作范围）：
% ---
% #### E.3.1 空间上的指数层与开集分类子
% #### E.3.2 群作用与 Grothendieck topos
% #### E.3.3 交换群值层与内射对象

\section{空间、群作用与 Grothendieck topos}
\label{b4:sec:E03}

预层范畴已经把函数和子对象都组织成了对象。加入层公理以后，局部相容性会改变其中的公式，却保留同样的泛性质。先在拓扑空间上把这件事算清楚，再比较群作用及一般位点。

\subsection{空间上的指数层与开集分类子}
\label{b4:subsec:E03:01}

\begin{proposition}\label{b4:prop:space-sheaf-topos}
设 \(X\) 为拓扑空间。集合值层范畴 \(\operatorname{Sh}_{\mathbf{Set}}(X)\) 是初等 topos。对层 \(F,G\)，指数层和子对象分类子分别为
\[
 (G^F)(U)=\operatorname{Hom}_{\operatorname{Sh}_{\mathbf{Set}}(U)}(F|_U,G|_U),
 \qquad
 \Omega(U)=\{\,V\subseteq U\mid V\text{ 开}\,\}.
\]
\(\Omega\) 的限制为 \(V\mapsto V\cap W\)，\(\mathrm{true}_U(*)=U\)。
\end{proposition}
\begin{proof}
终层、二元积和等化子逐开集构造后仍满足唯一拼接，故有限极限存在。对覆盖 \(U=\bigcup_iU_i\)，相容的局部层态射 \(F|_{U_i}\to G|_{U_i}\) 可逐截面拼接：对 \(s\in F(W)\)，先在 \(W\cap U_i\) 上求像，再用 \(G\) 的层公理拼成 \(W\) 上的截面。这些像与限制相容，且唯一。因此所给 \(G^F\) 是层。

求值在开集 \(U\) 上送 \((\theta,s)\) 到 \(\theta_U(s)\)。对层态射 \(b:H\times F\to G\)，\(z\in H(U)\) 决定局部态射
\[
 F(W)\longrightarrow G(W),\qquad
 s\longmapsto b_W(z|_W,s)\quad(W\subseteq U).
\]
这给出 \(H\to G^F\)，与求值构造互逆，所以为指数对象。

相容的局部开集由并集唯一拼接，故 \(\Omega\) 也是层。对子层 \(A\subseteq F\)，令
\[
 \chi_U(s)=\bigcup\{\,V\subseteq U\mid V\text{ 开且 }s|_V\in A(V)\,\}.
\]
这些局部截面都是 \(s\) 的限制，因而相容；\(A\) 的层公理使 \(s\) 在这整个并上属于 \(A\)。所以 \(\chi_U(s)\) 是使该限制属于 \(A\) 的最大开集。它与交 \(W\) 的限制相容，且等于 \(U\) 当且仅当 \(s\in A(U)\)，给出分类方块。若另一特征映射 \(\psi\) 也分类 \(A\)，则
\(W\subseteq\psi_U(s)\) 当且仅当 \(s|_W\in A(W)\)，故它必等于上述并集。这同时证明唯一性。
\end{proof}

\(\Omega(X)\) 是 \(X\) 的所有开集，不只是 \(\varnothing,X\)。在这个范畴中，“一个截面属于子层”可以在某个开区域成立。开集的补集未必开，因而这种子对象结构不总与普通集合的布尔代数相同。这里不需要另设逻辑公理，就能从局部限制看见差别。

\subsection{群作用与 Grothendieck topos}
\label{b4:subsec:E03:02}

\begin{example}\label{b4:exm:group-set-topos}
设 \(G\) 为群。左 \(G\)-集范畴的积为对角作用，指数对象的底集合为全部函数 \(X\to Y\)，作用为
\[
 (g\cdot f)(x)=g\cdot f(g^{-1}\cdot x).
\]
求值 \(Y^X\times X\to Y\) 等变；普通柯里化将等变映射 \(Z\times X\to Y\) 对应为等变映射 \(Z\to Y^X\)，直接核验上述公式即可。子对象是稳定子集，其特征函数到平凡作用的 \(\{0,1\}\) 等变，且唯一。因此该范畴为初等 topos。与预层的描述相比，群中每个箭头可逆，所以一个筛只有空或最大两种，恰好得到这两个真假值。
\end{example}

\begin{definition}\label{b4:def:grothendieck-topos}
设 \(\mathcal E\) 为局部小范畴。若存在小位点 \((\mathcal C,J)\)，使
\[
 \mathcal E\simeq\operatorname{Sh}_{\mathbf{Set}}(\mathcal C,J),
\]
则称 \(\mathcal E\) 为\term[Grothendieck-topos]{Grothendieck topos}[Grothendieck topos]。
\end{definition}

\begin{theorem}\label{b4:thm:site-topos}
设 \((\mathcal C,J)\) 为小位点。其集合值层范畴是初等 topos；因而每个 Grothendieck topos 都是初等 topos。
\end{theorem}
\begin{proofsketch}
记 \(\mathcal S=\operatorname{Sh}_{\mathbf{Set}}(\mathcal C,J)\)，\(i:\mathcal S\to\widehat{\mathcal C}\) 为包含函子。一般位点的层化函子 \(a:\widehat{\mathcal C}\to\mathcal S\) 满足 \(a\dashv i\)，并保持有限极限。其构造先以覆盖筛上的相容族定义 \(F^+\)，两族在共同覆盖细化上相同便视为等价，再作一次得到层 \(F^{++}\)。极限中的有限个相容条件可以在同一覆盖细化上核验，从而得到有限极限保持性。两次加号构造及这一步的细节见\cite[Chapter III, Section 5]{MacLaneMoerdijk1992Sheaves}。层的有限极限也可直接逐对象计算。

对层 \(F,G\)，先按\cref{b4:prop:presheaf-exponential}定义预层
\[
 E(U)\coloneqq\operatorname{Nat}(h_U\times iF,iG),
\]
限制沿 \(r:V\to U\) 由前复合 \(h_r\times1_{iF}\) 给出。一般位点上的 \(h_U\) 未必是层。由层化伴随、\(a\) 保持积及 \(a(iF)\cong F\)，有
\[
 E(U)\cong\operatorname{Hom}_{\mathcal S}(a(h_U)\times F,G).
\]
对任意层 \(H\)，同一伴随与 Yoneda 引理还给出
\[
 \operatorname{Hom}_{\mathcal S}(a(h_U),H)
 \cong\operatorname{Nat}(h_U,iH)\cong H(U).
\]

为核验 \(E\) 是层，取 \(R\in J(U)\) 及相容族 \(\theta_f\in E(V)\)，其中 \(f:V\to U\) 遍历 \(R\)。对每个 \(k:W\to V\)，相容性是
\(\theta_{f\circ k}=\theta_f\circ(h_k\times1_{iF})\)。对任意 \(g:V\to U\) 和 \(x\in F(V)\)，在覆盖筛 \(g^*R\) 上定义
\[
 y_k=(\theta_{g\circ k})_W(1_W,F(k)x)
 \quad(k:W\to V\text{ 属于 }g^*R).
\]
对 \(\ell:Z\to W\)，\(\theta_{g\circ k}\) 的自然性与上述相容性给出 \(G(\ell)y_k=y_{k\circ\ell}\)。因此 \(G\) 的层公理将这些截面唯一拼成 \(y\in G(V)\)，令 \(\theta_V(g,x)=y\)。沿 \(r:V'\to V\) 限制时，所用覆盖为 \(r^*(g^*R)=(g\circ r)^*R\)，局部值与同一公式相同；拼接的唯一性便给出 \(\theta\) 的自然性。当 \(g\in R\) 时，\(1_V\in g^*R\)，故
\(\theta_V(g,x)=(\theta_g)_V(1_V,x)\)；结合族的相容性，得到 \(E(f)\theta=\theta_f\)。任何另一拼接在每个 \(g^*R\) 上都具有相同的 \(y_k\)，由 \(G\) 的局部唯一性必等于 \(\theta\)。所以 \(E\) 是层。

现在对任意层 \(H\)，包含函子的全忠实性、预层指数的泛性质及层的积逐对象计算给出
\[
 \operatorname{Hom}_{\mathcal S}(H,E)
 \cong\operatorname{Nat}(iH\times iF,iG)
 \cong\operatorname{Hom}_{\mathcal S}(H\times F,G).
\]
这些双射关于 \(H\) 自然，求值由预层求值限制而来，故 \(E\) 是层范畴中的指数对象。

分类子取 \(\Omega(U)\) 为 \(U\) 上所有 \(J\)-闭筛组成的集合：筛 \(S\) 满足，对每个 \(f:V\to U\)，若 \(f^*S\in J(V)\)，则 \(f\in S\)。拉回保持这个条件，因为 \(h^*(f^*S)\) 覆盖时，\(S\) 的闭性给出 \(f\circ h\in S\)，即 \(h\in f^*S\)。因而筛的拉回定义预层 \(\Omega\)。

取 \(R\in J(U)\) 及相容族 \(S_f\in\Omega(V)\)，其中 \(f:V\to U\) 属于 \(R\)，相容性为 \(h^*S_f=S_{f\circ h}\)（\(h:W\to V\)）。先令
\[
 T=\{\,f:V\to U\mid f\in R,\ 1_V\in S_f\,\},
 \qquad
 S=\{\,g:V\to U\mid g^*T\in J(V)\,\}.
\]
若 \(f\in T\)，则 \(S_f\) 含恒等态射，因而是最大筛；于是 \(1_W\in h^*S_f=S_{f\circ h}\)，故 \(f\circ h\in T\)。这证明 \(T\) 对前复合封闭。对 \(f\in R\)，由相容性逐项得到
\[
 k\in f^*T
 \ \Longleftrightarrow\ 1_W\in S_{f\circ k}=k^*S_f
 \ \Longleftrightarrow\ k\in S_f
 \quad(k:W\to V),
\]
所以 \(f^*T=S_f\)。覆盖的拉回稳定性使 \(S\) 也对前复合封闭。若 \(q:V\to U\) 满足 \(q^*S\) 覆盖，则对其中每个 \(k\)，\(k^*(q^*T)=(q\circ k)^*T\) 都覆盖；拓扑的传递公理推出 \(q^*T\) 覆盖，故 \(q\in S\)。因此 \(S\) 是 \(J\)-闭筛。再对 \(f\in R\) 及 \(k:W\to V\)，有
\[
 k\in f^*S
 \ \Longleftrightarrow\ k^*(f^*T)=k^*S_f\in J(W)
 \ \Longleftrightarrow\ k\in S_f,
\]
最后一步使用 \(S_f\) 的闭性，以及筛含 \(k\) 时其沿 \(k\) 的拉回为最大筛。因此 \(f^*S=S_f\)，得到所需拼接。

若另一 \(J\)-闭筛 \(S'\) 也恢复该族，则 \(T\subseteq S'\)。对 \(g:V\to U\)，若 \(g\in S'\)，相容族的恢复给出 \(g^*T=g^*R\)，所以 \(g^*T\) 覆盖；反之，若 \(g^*T\) 覆盖，则 \(T\subseteq S'\) 使 \(g^*S'\) 覆盖，\(S'\) 的闭性给出 \(g\in S'\)。故 \(S'=S\)，证明 \(\Omega\) 的层条件。

取 \(\mathrm{true}_U(*)\) 为最大筛。上面的 \(\operatorname{Hom}_{\mathcal S}(a(h_U),H)\cong H(U)\) 说明单态射可逐对象识别为子层包含。对子层 \(A\subseteq F\)，令
\[
 \chi_U(x)=\{\,f:V\to U\mid F(f)x\in A(V)\,\}.
\]
若 \(f^*\chi_U(x)\) 覆盖，则局部限制在 \(A\) 中相容；由 \(A\) 的拼接与 \(F\) 的局部唯一性得到 \(F(f)x\in A(V)\)。所以 \(\chi_U(x)\) 是 \(J\)-闭筛。\(\chi\) 的自然性、其拉回恢复 \(A\) 及其唯一性，与\cref{b4:thm:presheaf-topos}中逐态射的核验相同，故 \(\Omega\) 是子对象分类子。这里省略的仍是一般筛上的两次加号构造及层化保持有限极限的技术细节；有关层化、指数和分类子的完整论证见\cite[Chapter III, Sections 4--7; Chapter IV, Section 2]{MacLaneMoerdijk1992Sheaves}。
\end{proofsketch}

初等 topos 的定义只规定有限极限、指数与分类子，不要求任意小余极限。例如有限集合范畴是初等 topos，却不是 Grothendieck topos：后者具有任意小余积，而前者没有可数个单元素集合的余积。由此可以区分两种名称的实际条件。

\subsection{交换群值层与内射对象}
\label{b4:subsec:E03:03}

集合值层可以带有内部的交换群结构：给出自然的加法 \(F\times F\to F\)、零截面及逆元，并逐开集满足群公理。它等价于通常的交换群值层。因此\secref{b4:sec:C03}采用的 \(\operatorname{Sh}(X)\) 是 \(\operatorname{Sh}_{\mathbf{Set}}(X)\) 中的交换群对象范畴，而不是这个 topos 本身。后者有函数对象与分类子，前者则有加法、核、余核与导出截面函子；各自的构造有不同用途。

\begin{theorem}\label{b4:thm:grothendieck-enough-injectives}
设 \(\mathcal A\) 为局部小交换范畴，具有小余积，滤过余极限存在且正合，并有生成元 \(U\)。则每个对象 \(M\) 都可单射到一个内射对象。
\end{theorem}
这就是\subsecref{b4:subsec:grothendieck-category}中 Grothendieck 范畴的内射存在性定理。原始论证见\cite[Théorème 1.10.1, pp. 135--137]{Grothendieck1957TohokuI}。
\begin{proofsketch}
生成元首先把子对象控制为集合。对任意对象 \(B\)，不同子对象 \(N\subseteq B\) 对应不同的子集 \(\operatorname{Hom}(U,N)\subseteq\operatorname{Hom}(U,B)\)，因为来自 \(U\) 的映射的像生成 \(N\)。因此 \(\operatorname{Sub}(B)\) 是集合；若 \(N\subseteq U\)，其子对象也都是 \(U\) 的子对象，故
\[
 |\operatorname{Sub}(N)|\le |\operatorname{Sub}(U)|.
\]
对每个 \(N\subseteq U\) 及态射 \(N\to M\)，取一份包含 \(N\to U\)。这些数据组成集合，可以将其直和沿所有 \(N\to M\) 给出的映射推出。直和包含是有限子直和的滤过余极限，因而由 AB5 仍为单态射。推出得到单态射 \(M\to M_1\)，使原有的每个 \(N\to M\) 都延到 \(U\to M_1\)。

选一个极限序数 \(\lambda\)，其共尾数满足
\[
 \operatorname{cf}(\lambda)>|\operatorname{Sub}(U)|.
\]
从 \(M_0=M\) 重复上述推出构造，在每个极限阶段取滤过余极限，得到连续单态射链 \((M_\alpha)_{\alpha\le\lambda}\)。AB5 保证各阶段到后继阶段及最终对象 \(I=M_\lambda\) 的映射仍单。

给定 \(N\subseteq U\) 及 \(f:N\to I\)，沿包含 \(\iota_\alpha:M_\alpha\to I\) 拉回，令 \(N_\alpha=f^{-1}(M_\alpha)\subseteq N\)。这些拉回也可写成
\[
 N_\alpha\cong\ker\bigl(N\oplus M_\alpha
 \xrightarrow{(f,-\iota_\alpha)}I\bigr).
\]
AB5 使取核与该滤过余极限相容，所得核为 \(\ker(N\oplus I\xrightarrow{(f,-1)}I)\cong N\)，因此
\[
 N=\bigcup_{\alpha<\lambda}N_\alpha.
\]
这个递增链只取 \(|\operatorname{Sub}(N)|\) 个可能的值。各个不同子对象首次出现的阶段组成大小小于 \(\operatorname{cf}(\lambda)\) 的集合，其上确界小于 \(\lambda\)；此后链已经稳定，故在某个 \(\alpha<\lambda\) 有 \(N_\alpha=N\)。这正表示 \(f\) 经 \(M_\alpha\) 因子化。由于 \(\lambda\) 为极限序数，\(\alpha+1<\lambda\)，下一次推出便将该态射延到 \(U\to M_{\alpha+1}\to I\)。

最后使用生成元版本的 Baer 判据：若 \(N\subseteq U\) 发出的每个态射都可延到 \(U\)，对象便内射。其理由是对任意 \(A\subseteq B\) 上的态射取极大延拓；\(\operatorname{Sub}(B)\) 是集合，局部小性又使全部延拓组成集合，而滤过余极限正合使延拓的链有上界。若延拓的定义域尚未达到 \(B\)，生成元提供一个像不包含于该定义域的 \(U\to B\)，拉回得到 \(N\subseteq U\)，再经推出扩张，矛盾。于是 \(I\) 内射。这里给出了统一阶段界的机制；超限递归及生成元判据的原始论证见所引原文。
\end{proofsketch}

模范畴中可以直接构造内射模；拓扑空间上的交换群值层也在\cref{b4:prop:C-sheaf-injectives}中通过茎得到内射嵌入。这个一般定理说明，内射消解的可用性还可以由范畴的正合余极限与生成元保证，而不必依赖每一类对象的元素模型。

% !TeX root = ../../Book4.tex
% 纲要标题：### E.4 平展位点与平展 topos
% 纲要细目（写作范围）：
% #### E.4.1 平展覆盖与层
% #### E.4.2 域上的平展层与 Galois 作用
% ---

\section{平展位点与平展 topos}
\label{b4:sec:E04}

在 \(\operatorname{Spec}\mathbb R\) 上，方程 \(T^2+1=0\) 不能通过缩小开集得到解，却能在平展态射 \(\operatorname{Spec}\mathbb C\rightarrow\operatorname{Spec}\mathbb R\) 后得到解。这样的局部变化仍有交叠与拼接，只是交叠须用纤维积表示。第三卷第24.4节已将环的平展性推广到概形态射，并证明它保持于复合与基变换；这里用这些态射构造位点，再将\subsecref{b4:subsec:C04:02}的域上计算放到层范畴中。

\subsection{平展覆盖与层}
\label{b4:subsec:E04:01}

先固定一个含概形 \(X\) 的集合大小的概形框架：其中的概形及全部态射组成小范畴，并且开子概形、纤维积及框架中概形上的有限型概形可在其中选取同构代表。为比较不同框架，先对 \(X\) 的每个仿射开集 \(W=\operatorname{Spec}A\)，从有限展示平展 \(A\) 代数的各同构类中选取代表 \(B\)，保留结构态射 \(\operatorname{Spec}B\rightarrow W\rightarrow X\)。这些局部模型只有集合多个：\(X\) 的仿射开集构成集合，而有限组生成元与关系所给的 \(A\) 代数只有集合多个展示。要求所选框架含每个局部模型的一个 \(X\)-同构代表；继续加入所需的开集、有限型概形与纤维积，仍可在一个足够大的集合中完成这些选择。以下构造都在这个框架内进行，取同构代表时同时保留结构态射。

\begin{definition}\label{b4:def:small-etale-site}
设 \(X\) 为概形，在上述集合框架中取范畴 \(\mathcal C_X\)：其对象为平展概形态射 \(p_U:U\rightarrow X\)，从 \(U\) 到 \(V\) 的态射为满足 \(p_V\circ f=p_U\) 的概形态射 \(f:U\rightarrow V\)。对对象 \(U\)，若一族态射 \((u_i:U_i\rightarrow U)_{i\in I}\) 中每个 \(u_i\) 平展，且
\[
 U=\bigcup_{i\in I}u_i(U_i),
\]
则称其为 \(U\) 的\term[pingzhan-fugai]{平展覆盖}[étale covering]。规定一个筛覆盖 \(U\)，当且仅当它包含一族这样的覆盖态射，所得位点记为 \(X_{\acute{e}t}\)，称为 \(X\) 的\term[xiao-pingzhan-weidian]{小平展位点}[small étale site]。
\end{definition}

名称中的“小”是相对于允许一般 \(X\)-概形作为对象的大平展位点而言的；对象的结构态射须平展。范畴在集合意义下为小，是由前面的框架选择保证的。任意大的不交并并不必全都属于这个已选范畴。另一方面，每个平展对象都可由映入 \(X\) 的仿射开集的仿射平展模型覆盖，这些模型来自有限展示，已经可以在框架内选到。

\begin{proposition}\label{b4:prop:etale-covering-topology}
设 \(X\) 为概形，\(\mathcal C_X\) 为\cref{b4:def:small-etale-site}中的小范畴。平展覆盖满足恒等覆盖、复合覆盖和基变换覆盖条件；定义中的覆盖筛给出\cref{b4:def:grothendieck-site}的 Grothendieck 拓扑。
\end{proposition}
\begin{proof}
恒等态射平展且满射。若 \((U_i\rightarrow U)\) 覆盖 \(U\)，各 \((U_{ij}\rightarrow U_i)\) 又覆盖 \(U_i\)，则复合态射平展，其像的并仍为 \(U\)。

对任意 \(f:V\rightarrow U\)，基变换 \(V\times_UU_i\rightarrow V\) 平展，源也平展于 \(X\)，故仍是 \(\mathcal C_X\) 中的对象。它们的像覆盖 \(V\)：对 \(v\in V\)，取 \(u_i\in U_i\) 映到 \(f(v)=u\)；非零环
\[
 \kappa(v)\otimes_{\kappa(u)}\kappa(u_i)
\]
有素理想，给出纤维积中映到 \(v\) 的点。这里张量积非零，是因为两个因子都是 \(\kappa(u)\) 上的非零向量空间。

最大筛含恒等覆盖。若筛含 \((u_i)\)，它沿 \(f\) 的拉回便含刚构造的基变换覆盖。最后，设 \(S\) 覆盖 \(U\)，而筛 \(R\) 对每个 \(s\in S\) 的拉回都覆盖。选 \(S\) 中一族覆盖 \((u_i)\)，再在各 \(u_i^*R\) 中选覆盖 \((v_{ij})\)。复合 \((u_i\circ v_{ij})\) 覆盖 \(U\)，且全属于 \(R\)，故 \(R\) 覆盖。这正是三条筛公理。
\end{proof}

\begin{proposition}\label{b4:prop:etale-sheaf-equalizer}
设 \(X\) 为概形，\(F:\mathcal C_X^{\mathrm{op}}\rightarrow\mathbf{Set}\) 为预层。则 \(F\) 是 \(X_{\acute{e}t}\) 上的层，当且仅当对每个平展覆盖 \((U_i\rightarrow U)\)，下列图是等化子图：
\begin{equation}\label{b4:eq:etale-sheaf-equalizer}
 F(U)\longrightarrow\prod_iF(U_i)
 \mathrel{\substack{\longrightarrow\\[-0.6ex]\longrightarrow}}
 \prod_{i,j}F(U_i\times_UU_j).
\end{equation}
右侧两箭头分别由两投影的限制给出。
\end{proposition}
\begin{proof}
设 \(S\) 为这族覆盖生成的筛。一个相容族 \(x_i\in F(U_i)\) 对每个分解 \(f=u_i\circ g:V\rightarrow U\) 给出 \(F(g)(x_i)\)。若还有 \(f=u_j\circ h\)，则 \((g,h)\) 给出到 \(U_i\times_UU_j\) 的态射，等化子中的相容条件保证两值相同。这些值与进一步限制相容，因而给出 \(S\rightarrow F\) 的自然变换。反过来，取该变换在各 \(u_i\) 上的值，便恢复相容族。于是\eqref{b4:eq:etale-sheaf-equalizer}恰是生成筛上的层条件。

若任意覆盖族都满足这一条件，给定任意覆盖筛 \(R\) 上的相容族，可先选 \(R\) 中一族覆盖，将其值唯一拼成 \(x\in F(U)\)。对任意 \(f:V\rightarrow U\) 属于 \(R\)，沿拉回覆盖检验，\(F(f)(x)\) 与原相容族在 \(f\) 上的值相同；等化子的单射性使它们在 \(V\) 上相同。故 \(x\) 恢复全部相容族，并且唯一。反方向由生成筛是覆盖筛立即得到。
\end{proof}

\begin{proposition}\label{b4:prop:etale-affine-basis}
设 \(X\) 为概形。在 \(X_{\acute{e}t}\) 中只取仿射源的对象，并以仿射对象组成的联合满射平展族为覆盖，所得位点的层范畴与 \(X_{\acute{e}t}\) 的层范畴等价。
\end{proposition}
\begin{proof}
每个对象 \(U\) 有仿射开覆盖 \((U_\lambda)\)。给定仿射子位点上的层 \(F\)，定义 \(U\) 上的截面为 \(x_\lambda\in F(U_\lambda)\) 组成的相容族；相容是指在每个 \(U_\lambda\cap U_\mu\) 的仿射开覆盖上，两者限制相同。两组仿射开覆盖的交可以再作仿射开覆盖，仿射对象上的层公理使相容族沿细化的映射为双射。因此定义与所选开覆盖无关。

态射的拉回在仿射开集的原像上再取仿射开覆盖即可定义，细化无关性保证复合律。对任意平展覆盖，其在这些仿射开集上的基变换可再被仿射开集覆盖；局部层公理先在这些仿射对象上拼接，再拼接到 \(U\)，证明所构造的预层是层。限制后恢复 \(F\)，而原位点上的层由其仿射开覆盖上的值唯一恢复；层态射也由这些限制唯一恢复。这给出互逆的函子。

还可用同一拼接论证比较两个允许的集合框架。记节首选取的局部模型组成的小范畴为 \(\mathcal B_X\)，其态射取全部 \(X\)-态射，覆盖仍为联合满射平展族。每个框架中的模型代表经 \(X\)-同构与 \(\mathcal B_X\) 对应；共轭这些同构给出态射的对应，并保持覆盖，因为同构不改变态射像的并。

对任意平展对象 \(U\rightarrow X\)，先将 \(U\) 按 \(X\) 的仿射开集 \(W\) 的原像覆盖，再在各原像中取仿射开覆盖。所得仿射对象平展于 \(W\)，其坐标代数有限展示，故都与 \(\mathcal B_X\) 中的模型 \(X\)-同构。两个这样的覆盖的交叠，以及沿任意态射拉回所得的对象，都能再按同样方式被这些局部模型覆盖。因此刚才的相容族、细化与拉回构造可全部在 \(\mathcal B_X\) 上进行，并从它恢复每个框架中的层及层态射。两项限制函子
\[
 \operatorname{Sh}_{\mathbf{Set}}(\mathcal C_X,J)
 \longrightarrow\operatorname{Sh}_{\mathbf{Set}}(\mathcal B_X)
\]
遂分别为等价，给出两个框架所得层范畴的自然等价；改变代表或所选覆盖，只改变恢复函子的自然同构。这里共同的是一组覆盖每个对象的局部模型，并非全部仿射源对象的同构类型：仿射源对象的像未必落在 \(X\) 的一个仿射开集中。
\end{proof}

两个仿射对象在非仿射对象上的纤维积未必仿射。使用仿射基时，\eqref{b4:eq:etale-sheaf-equalizer}中的相容性便须在这个纤维积的仿射开覆盖上检验，而不能把它直接当作仿射基中的一个对象。

\begin{definition}\label{b4:def:etale-topos}
设 \(X\) 为概形。小平展位点的集合值层范畴
\[
 \operatorname{Sh}_{\mathbf{Set}}(X_{\acute{e}t})
\]
称为 \(X\) 的\term[pingzhan-topos]{平展 topos}[étale topos]。其交换群对象组成的范畴记为 \(\operatorname{Sh}_{\mathbf{Ab}}(X_{\acute{e}t})\)；对象等价地是取值于交换群并满足\eqref{b4:eq:etale-sheaf-equalizer}的预层。
\end{definition}

由\cref{b4:thm:site-topos}，集合值平展层构成 Grothendieck topos，具有有限极限、指数对象及子对象分类子。交换群平展层则有加法与正合列；对其取全局截面的导出函子，需要内射消解。两种范畴的联系由\subsecref{b4:subsec:E03:03}中的内部交换群结构给出。

\begin{proposition}\label{b4:prop:abelian-site-sheaves-grothendieck}
设 \((\mathcal C,J)\) 为小位点。交换群值层范畴 \(\operatorname{Sh}_{\mathbf{Ab}}(\mathcal C,J)\) 是 Grothendieck 范畴，因而有足够多内射对象。
\end{proposition}
\begin{proofsketch}
在\cref{b4:thm:site-topos}的两次加号构造中，相容族逐项相加，得到加法层化函子 \(a\)。它是交换群值层包含到预层范畴的左伴随，并保持有限极限。因此它保持核；作为加性左伴随又保持余核，故正合。层的核逐对象计算，余核为逐对象余核的层化，由此得到交换范畴。

任意小余极限都是先在预层中逐对象取余极限，再层化。交换群的滤过余极限与核交换，而层的核也逐对象计算；再用正合层化，便知层范畴的滤过余极限保持核。余极限作为左伴随保持余核，因此这些滤过余极限正合。对可表预层 \(h_U=\operatorname{Hom}_{\mathcal C}(-,U)\)，记 \(\mathbb Z[h_U]\) 为逐对象在这个集合上取自由交换群的预层。伴随与 Yoneda 引理给出
\[
 \operatorname{Hom}\bigl(a\mathbb Z[h_U],F\bigr)\cong F(U).
\]
由于 \(\mathcal C\) 小，可取层的余积 \(G=\bigoplus_{U\in\mathcal C}a\mathbb Z[h_U]\)。若两个层态射不同，它们在某个 \(F(U)\) 的元素上不同，相应的 \(a\mathbb Z[h_U]\rightarrow F\) 检测这个差别；令其余直和项映为零，就得到检测差别的 \(G\rightarrow F\)。故 \(G\) 是生成元。这里使用的层化伴随与有限极限保持性，正是\cref{b4:thm:site-topos}中给出的层化概要；其余各项核验说明可以应用\cref{b4:thm:grothendieck-enough-injectives}。
\end{proofsketch}

在 \(X_{\acute{e}t}\) 中，\(X\rightarrow X\) 是终对象。对交换群平展层 \(F\)，全局截面函子 \(\Gamma(X,F)=F(X)\) 左正合，因而定义
\[
 H^n_{\acute{e}t}(X,F)=R^n\Gamma(X,-)(F)
 \qquad(n\ge0).
\]
这给出\secref{b4:sec:C04}所用平展上同调的位点定义。

\subsection{域上的平展层与 Galois 作用}
\label{b4:subsec:E04:02}

\begin{definition}\label{b4:def:finite-etale-object}
设 \(X\) 为概形，\(p:U\rightarrow X\) 为平展态射。若对每个仿射开集 \(\operatorname{Spec}A\subseteq X\)，原像都为仿射概形 \(\operatorname{Spec}B\)，且 \(B\) 作为 \(A\) 模有限生成，则称 \(p\) 为有限平展态射，\(U\) 为有限平展 \(X\)-对象。
\end{definition}

一般平展对象不能只用有限平展对象代替。例如取域 \(k\)，令 \(X=\operatorname{Spec}k[t]\)、\(U=D(t)\)。开浸入 \(U\rightarrow X\) 平展；但若非空有限平展 \(V\rightarrow X\) 能经过 \(U\) 分解，写 \(V=\operatorname{Spec}B\)，则 \(B\) 是非零有限平坦 \(k[t]\) 代数。它作为主理想整环上的有限生成无挠模为正秩自由模，故 \(B/tB\ne0\)；分解到 \(D(t)\) 却要求 \(t\) 在 \(B\) 中可逆，矛盾。在域上，所需的有限性由下面的论证产生。

\begin{lemma}\label{b4:lem:etale-over-field-components}
设 \(K\) 为域，\(U\rightarrow\operatorname{Spec}K\) 为平展概形态射。则存在一族有限可分域扩张 \((L_\lambda/K)_{\lambda\in\Lambda}\)，使
\[
 U\cong\coprod_{\lambda\in\Lambda}\operatorname{Spec}L_\lambda.
\]
反过来，这样的不交并平展于 \(K\)。特别地，仿射平展 \(K\)-概形对应有限个有限可分域的直积。
\end{lemma}
\begin{proof}
先设 \(B\) 为平展 \(K\) 代数；零代数对应空对象，以下设 \(B\ne0\)。它有限展示，且 \(\Omega_{B/K}=0\)，也可由第三卷第23.3节的标准平展模型刻画。取代数闭包 \(\overline K\)，令 \(C=B\otimes_K\overline K\)。基变换使 \(C\) 仍有限展示且 \(\Omega_{C/\overline K}=0\)。由第三卷第6.2节的弱零点定理，每个极大理想 \(\mathfrak m\) 的剩余域为 \(\overline K\)。第三卷第14.3节的有理点余切空间同构给出
\[
 \mathfrak mC_{\mathfrak m}/(\mathfrak mC_{\mathfrak m})^2
 \cong\overline K\otimes_C\Omega_{C/\overline K}=0.
\]
\(C_{\mathfrak m}\) 是 Noether 局部环，Nakayama 引理使其极大理想为零，所以它是域。每个素理想包含于某个极大理想；局部化中的素理想对应说明这两者必相等。幂零元在每个极大理想处局部化后也为零，故本身为零。因此 \(C\) 约化且零维，由第三卷第3.3节的 Noether 零维环定理，\(C\) 是有限个 \(\overline K\) 的直积。

\(B\rightarrow C\) 单射，故 \(B\) 约化。\(K\) 上线性无关的元素在扩域后仍线性无关，而 \(C\) 是有限维 \(\overline K\) 向量空间，所以 \(B\) 是有限维 \(K\) 代数。约化 Artin 环的分解使 \(B\) 为有限个有限扩张域的直积。微分模在各因子上仍为零，第三卷第14.4--14.5节的有限域扩张判据说明每个因子可分。

对一般 \(U\)，平展定义给每一点一个上述仿射平展邻域。这个邻域是有限个域的谱的不交并，因而每一点本身就是开子概形，且其环为相应的有限可分域。将全部这些单点开子概形取不交并即恢复 \(U\)。反方向逐分量由有限可分域扩张平展的判据得到。若 \(U\) 仿射，它准紧，其单点开覆盖必有有限子覆盖，故只有有限个分量。
\end{proof}

\begin{proposition}\label{b4:prop:field-finite-etale-basis}
设 \(K\) 为域。有限平展 \(K\)-概形在 \((\operatorname{Spec}K)_{\acute{e}t}\) 中构成覆盖基：每个对象都有来自这些对象的覆盖。取其同构代表组成小范畴 \(\mathrm{F\acute Et}(K)\)，覆盖仍为联合满射族，则限制给出集合值层范畴的等价。这个等价也适用于交换群值层。
\end{proposition}
\begin{proof}
前一引理的分量 \(U_\lambda=\operatorname{Spec}L_\lambda\) 就给出所需覆盖。有限平展对象的纤维积仍有限平展：平展性保持基变换与复合，其坐标代数的张量积仍有限维。有限可分域扩张可由多项式的根表示，只有集合多个同构类型，故其有限直积也可选小范畴中的代表。

给定有限平展基上的层 \(F\)，对任意 \(U=\coprod_\lambda U_\lambda\) 定义
\[
 \widetilde F(U)=\prod_\lambda F(U_\lambda).
\]
对态射 \(f:V\rightarrow U\)，每个分量 \(V_\mu\) 映到唯一 \(U_{\lambda(\mu)}\)；沿这个分量态射限制相应坐标，定义 \(\widetilde F(f)\)。分量态射的复合律保证这是预层。基中的有限不交并由分量覆盖，分量间交为空，且空覆盖给出 \(F(\varnothing)\) 是单元素集合；故 \(\widetilde F\) 限制后自然恢复 \(F\)。

设 \((V_i\rightarrow U)\) 为平展覆盖，并给定相容的截面。固定一个 \(U_\lambda\)，选 \(V_i\) 中一个映到它的分量 \(W\)。\(W\rightarrow U_\lambda\) 是非空有限可分扩张给出的覆盖。该分量截面在 \(W\times_{U_\lambda}W\) 上相容，由 \(F\) 的层公理唯一下降到 \(x_\lambda\in F(U_\lambda)\)。对另一个映到 \(U_\lambda\) 的分量 \(W'\)，两项在 \(W\times_{U_\lambda}W'\) 上相同；这个纤维积覆盖 \(W'\)，故层的单射性使 \(x_\lambda\) 的限制恢复 \(W'\) 上的截面。所有 \(x_\lambda\) 因而唯一拼成 \(\widetilde F(U)\) 的元素，证明层条件。

原位点上的任意层也由分量开覆盖给出 \(F(U)\cong\prod_\lambda F(U_\lambda)\)，所以扩展唯一，层态射逐分量恢复。这证明等价。取值于交换群时，同一构造逐项保留加法，空对象取零群，论证不变。
\end{proof}

固定可分闭包 \(K^s\)，记 \(G_K=\operatorname{Gal}(K^s/K)\)。\subsecref{b4:subsec:C04:02}已经将有限平展 \(K\)-概形 \(U\) 对应到有限连续 \(G_K\)-集
\[
 T(U)=\operatorname{Hom}_K(\operatorname{Spec}K^s,U),
\]
在 \(U=\operatorname{Spec}B\) 时，几何点就是 \(K\)-代数同态 \(t:B\rightarrow K^s\)，作用为 \(g\cdot t=g\circ t\)。该小节通过 \(G_K\)-等变函数 \(T(U)\rightarrow K^s\) 恢复 \(B\)，并以函数的拉回恢复代数同态；与 \(\operatorname{Spec}\) 的反变方向相抵，得到自然双射
\[
 \operatorname{Hom}_K(U,V)
 \cong\operatorname{Hom}_{G_K}(T(U),T(V)),
 \qquad f\longmapsto T(f),
\]
其中 \(T(f)(t)=f\circ t\)。任意有限连续 \(G_K\)-集是有限个 \(G_K/H\) 的不交并，开子群 \(H\) 对应有限可分域 \((K^s)^H\)，所以这一全忠实函子也是本质满的。它将覆盖对应到联合满射。

\ProseParagraphBreak
对任意离散连续 \(G_K\)-集 \(M\)，在有限平展基上定义
\[
\begin{aligned}
 F_M(U)&=\operatorname{Hom}_{G_K}(T(U),M),\\
 F_M(f)(h)&=h\circ T(f)\qquad(f:V\rightarrow U).
\end{aligned}
\]
联合满射覆盖上的相容等变函数唯一拼接，故 \(F_M\) 是层，再由\cref{b4:prop:field-finite-etale-basis}扩展到整个小平展位点。这里 \(M\) 可以无限，连续性要求每个元素的稳定子群开。

\ProseParagraphBreak
反过来，给定集合值平展层 \(F\)，简记 \(F(L)=F(\operatorname{Spec}L)\)，沿 \(K^s/K\) 中的有限 Galois 子扩张取
\[
 M_F=\varinjlim_{L/K\ {\rm finite\ Galois}}F(L).
\]
沿域包含的过渡映射为拉回；相应概形态射是覆盖，故层的分离性使这些映射单射。对 \(g\in G_K\)，在 \(F(L)\) 上令 \(g\) 通过代数自同构 \(g|_L\) 的拉回 \(F(\operatorname{Spec}(g|_L))\) 作用。\(\operatorname{Spec}\) 与 \(F\) 的两次反变使 \(g\cdot(h\cdot s)=(gh)\cdot s\)，得到左作用；域自同构与包含相容，故这些作用也与过渡映射相容。每个来自 \(F(L)\) 的元素被 \(G_L=\operatorname{Gal}(K^s/L)\) 固定，所以 \(M_F\) 是离散连续 \(G_K\)-集。

\ProseParagraphBreak
对有限平展 \(U=\operatorname{Spec}B\)，构造
\[
 \theta_U:F(U)\longrightarrow
 \operatorname{Hom}_{G_K}(T(U),M_F).
\]
若 \(s\in F(U)\)、\(t:B\rightarrow K^s\)，取有限 Galois 子扩张 \(L/K\) 包含 \(t(B)\)，将 \(s\) 沿 \(\operatorname{Spec}L\rightarrow U\) 拉回，所得截面在余极限中的类就是 \(\theta_U(s)(t)\)。增大 \(L\) 不改变这个类；拉回复合律还给出
\[
 \theta_U(s)(g\cdot t)=g\cdot\theta_U(s)(t).
\]
因此这确实是按几何点取茎得到的等变函数。

\ProseParagraphBreak
选有限 Galois \(L/K\) 使 \(U\) 在 \(L\) 上分裂，则
\(U\times_K\operatorname{Spec}L\cong\coprod_{t\in T(U)}\operatorname{Spec}L\) 覆盖 \(U\)。若两个截面经 \(\theta_U\) 得到同一函数，它们在每个上述分量上的拉回在 \(M_F\) 中相同。由 \(F(L)\rightarrow M_F\) 单射，这些拉回已经相同，再由层的分离性得到原截面相同，故 \(\theta_U\) 单射。

\ProseParagraphBreak
给定等变函数 \(h:T(U)\rightarrow M_F\)，由于 \(T(U)\) 有限，可增大分裂域 \(L\)，使所有 \(h(t)\) 都由 \(x_t\in F(L)\) 表示。\(F(L)\rightarrow M_F\) 单射和 \(h\) 的等变性给出
\[
 x_{\sigma\cdot t}=\sigma\cdot x_t
 \qquad(\sigma\in\operatorname{Gal}(L/K)).
\]
按 \(a\otimes b\mapsto(a\sigma(b))_\sigma\) 识别 \(L\otimes_KL\cong\prod_{\sigma\in\operatorname{Gal}(L/K)}L\)，分裂覆盖的交叠在 \((t,\sigma)\) 分量上的两投影给出 \(x_t\) 与 \(\sigma\cdot x_{\sigma^{-1}\cdot t}\)，上述等式使它们相同。层公理使 \((x_t)_t\) 唯一下降为 \(s\in F(U)\)，且 \(\theta_U(s)=h\)。这证明满射，因而 \(\theta_U\) 为双射。

\ProseParagraphBreak
这项双射与全部限制相容。对任意有限平展对象间的态射 \(f:V\rightarrow U\)，拉回的复合律给出
\[
 \theta_V\bigl(F(f)(s)\bigr)(t)
 =\theta_U(s)\bigl(T(f)(t)\bigr)
 \qquad(t\in T(V)).
\]
故 \(\theta:F\xrightarrow{\sim}F_{M_F}\) 在有限平展基上是层的自然同构，再由基的恢复等价扩展为整个小平展位点上的自然同构。

\ProseParagraphBreak
层态射 \(\eta:F\rightarrow F'\) 在各 \(F(L)\) 上给出相容映射，取余极限得到等变映射 \(M_F\rightarrow M_{F'}\)。反过来，等变映射 \(a:M\rightarrow M'\) 通过 \(h\mapsto a\circ h\) 给出层态射 \(F_M\rightarrow F_{M'}\)。取茎的定义使 \(\theta\) 与 \(\eta\) 相容，因此这两个态射构造相互恢复。从 \(M\) 出发时，在单位嵌入 \(L\hookrightarrow K^s\) 处求值给出
\[
 F_M(\operatorname{Spec}L)\cong M^{G_L},
 \qquad
 M_{F_M}\cong\varinjlim_L M^{G_L}=M.
\]
最后的等式成立，因为每个元素的开稳定子群包含某个 \(G_L\)；该同构与等变映射相容。它和 \(\theta\) 给出两复合函子的自然同构，证明集合值平展层范畴与离散连续 \(G_K\)-集范畴等价。

\ProseParagraphBreak
终对象 \(\operatorname{Spec}K\) 对应单元素的平凡 \(G_K\)-集，故 \(\Gamma(\operatorname{Spec}K,F_M)=M^{G_K}\)。对交换群对象，所得等价是 C.4.2 的交换群平展层与离散连续 \(G_K\)-模的正合等价；全局截面对应不变量，导出后便得到\eqref{b4:eq:C-etale-field}。有限的是作为覆盖基的几何对象，层的截面与相应的 \(G_K\)-模均不要求有限。

\begin{example}\label{b4:exm:real-etale-topos}
取 \(X=\operatorname{Spec}\mathbb R\)。其 Zariski 底空间只有一个点，集合值 Zariski 层范畴等价于 \(\mathbf{Set}\)。可分闭包为 \(\mathbb C\)，绝对 Galois 群为 \(C_2=\{1,\sigma\}\)，其中 \(\sigma\) 是复共轭；平展 topos 则等价于带有 \(C_2\) 作用的集合范畴，有限群上的这些离散作用都连续。

例如令 \(M=\{+,-\}\)，\(\sigma\) 交换两个元素。对应的层满足
\[
 F_M(\operatorname{Spec}\mathbb C)\cong M,
 \qquad F_M(X)=M^{C_2}=\varnothing.
\]
局部截面存在，却没有不变的整体截面。用
\[
 \mathbb C\otimes_{\mathbb R}\mathbb C\xrightarrow{\sim}
 \mathbb C\times\mathbb C,\qquad
 a\otimes b\longmapsto(ab,a\overline b)
\]
识别交叠后，层条件中的两箭头为 \(m\mapsto(m,m)\) 和 \(m\mapsto(m,\sigma m)\)，等化子正是固定点集。对离散 \(C_2\)-模 \(A\)，\eqref{b4:eq:C-etale-field}则给出每个次数的识别
\[
 H^n_{\acute{e}t}(\operatorname{Spec}\mathbb R,F_A)
 \cong H^n(C_2,A)\qquad(n\ge0).
\]
集合值层的作用与固定点描述因而也决定了交换群层所用的系数范畴及导出截面计算。
\end{example}


\begin{chaptersummary}
筛把一组覆盖箭头及其全部限制组织为子预层；层要求筛上的相容数据来自唯一截面。指数对象记录与所有限制相容的映射，分类子记录一个截面在哪些限制下属于子对象。对于预层，这些限制组成筛；对于空间上的层，它们组成开集。

\ProseParagraphBreak
初等 topos 用有限极限、指数与分类子刻画这些结构；Grothendieck topos 进一步要求来自一个小景的集合值层。交换群值层则是其中的交换群对象，具有用于同调的正合结构。对它们计算导出截面时，内射消解的存在另由具体构造或 Grothendieck 范畴定理保证。

\ProseParagraphBreak
平展景中的局部对象可以改变剩余域，交叠由纤维积给出，层仍要求相容截面唯一拼接。域上的平展对象由有限可分扩张的谱组成，有限平展对象因而足以计算层。集合值层对应离散连续 Galois 集，交换群值层对应离散连续 Galois 模；全局截面取固定元素，导出截面便成为连续 Galois 上同调。
\end{chaptersummary}

\BookExercises

\begin{exercise}\label{b4:ex:E-sieve-downset}
把偏序集看成范畴。证明对象 \(u\) 上的筛恰是 \(\{v:v\le u\}\) 中的向下封闭子集。列出链 \(0<1<2\) 中对象 \(2\) 上的全部筛。
\end{exercise}
\begin{solution}
指向 \(u\) 的箭头由 \(v\le u\) 唯一确定；前复合就是进一步选择 \(w\le v\)。所以筛条件恰好是向下封闭。所求四个筛为 \(\varnothing,\{0\},\{0,1\},\{0,1,2\}\)。
\end{solution}

\begin{exercise}\label{b4:ex:E-empty-cover}
若景上某对象 \(U\) 的空筛是覆盖，证明每个集合值层都有 \(F(U)\) 为单元素集合。解释开集景中空集的情形。
\end{exercise}
\begin{solution}
空预层到 \(F\) 的自然变换唯一，所以 \(\operatorname{Nat}(\varnothing,F)\) 为单元素集合。层条件要求 \(F(U)\) 与它双射。空集由空覆盖覆盖，故 \(F(\varnothing)\) 应是单元素集合；交换群值时这就是零群的底集合。
\end{solution}

\begin{exercise}\label{b4:ex:E-powerset}
设 \(G\) 为群，\(X\) 为左 \(G\)-集。将 \(G\)-集范畴中的幂对象 \(\Omega^X\) 的底集合识别为 \(X\) 的幂集，写出其 \(G\)-作用，并说明固定元素对应哪些子集。对 \(G=C_2\) 交换两点的传递集 \(X\)，分别计算幂对象的元素数与全局元素数。
\end{exercise}
\begin{solution}
\(\Omega=\{0,1\}\) 带平凡作用。函数 \(f:X\to\Omega\) 与子集 \(A=f^{-1}(1)\) 互相确定。指数上的作用为 \((g\cdot f)(x)=f(g^{-1}\cdot x)\)，所以 \((g\cdot f)^{-1}(1)=gA\)。因此幂集上的作用是 \(A\mapsto gA\)，固定元素恰是 \(G\)-稳定子集；这正好与 \(X\) 在 \(G\)-集范畴中的子对象对应。

写 \(X=\{x_0,x_1\}\)，其幂集的四个元素为 \(\varnothing,\{x_0\},\{x_1\},X\)。非恒等群元素交换两个单元素子集，并固定空集与全集。因而幂对象有四个底集合元素，只有两个全局元素：从终对象到幂对象的等变映射必须取固定元素。
\end{solution}

\begin{exercise}\label{b4:ex:E-group-exponential}
设 \(G\) 为群，\(X,Y\) 为左 \(G\)-集。证明指数 \(Y^X\) 的固定元素恰为等变映射 \(X\to Y\)，并解释为什么底集合应取全部函数。
\end{exercise}
\begin{solution}
固定条件为 \(g\cdot f(g^{-1}\cdot x)=f(x)\)，令 \(x=g\cdot z\) 便得 \(f(g\cdot z)=g\cdot f(z)\)。反向代入也成立。指数的底集合需允许每个函数，以承载其共轭作用；只取等变函数会只留下固定部分，无法将任意等变 \(Z\times X\to Y\) 柯里化为所需的带作用对象。
\end{solution}

\begin{exercise}\label{b4:ex:E-classifier-zero-locus}
令 \(F\) 为实直线上的连续实值函数层，\(A\) 为零函数子层。计算 \(\chi_U(f)\)，并对 \(U=\mathbb R\)、\(f(x)=\max(x,0)\) 给出具体结果。
\end{exercise}
\begin{solution}
\(\chi_U(f)\) 是 \(f^{-1}(0)\) 在 \(U\) 中的内部，因为它是使限制恒零的最大开集。给定函数的零点集为 \(({-\infty},0]\)，其内部为 \(({-\infty},0)\)。分类子因此记录开区域，不是闭零点集本身。
\end{solution}

\begin{exercise}\label{b4:ex:E-finite-sets}
证明有限集合范畴具有有限极限、指数与子对象分类子，却不存在可数个单元素集合的余积。
\end{exercise}
\begin{solution}
有限极限由有限积的子集构造，仍有限；有限集间的全部函数构成有限指数；\(\{0,1\}\) 仍分类全部子集。若所说余积为有限集 \(Q\)，设各结构映射的像为 \(q_n\)。任意到 \(\{0,1\}\) 的取值族必须由函数 \(Q\to\{0,1\}\) 实现，因此 \(q_n\) 两两不同，否则在相同的两个位置指定不同值便矛盾。有限 \(Q\) 无法含这些点。
\end{solution}

\begin{exercise}\label{b4:ex:E-abelian-topos}
证明具有零对象的笛卡尔闭范畴中，每个对象都同构于零对象。由此说明非零模范畴为什么不是初等 topos。
\end{exercise}
\begin{solution}
笛卡尔闭性使 \(X\times-\) 有右伴随，因而保持始对象。零对象 \(0\) 也是终对象，所以一方面 \(X\times0\cong0\)，另一方面 \(X\times0\cong X\)，故 \(X\cong0\)。非零环上的模范畴有零对象及非零对象，不可能笛卡尔闭，因而不是初等 topos。
\end{solution}

\begin{exercise}\label{b4:ex:E-real-etale-overlap}
证明映射
\[
\mathbb C\otimes_{\mathbb R}\mathbb C\longrightarrow\mathbb C\times\mathbb C,
\qquad z\otimes w\longmapsto(zw,z\overline w)
\]
是 \(\mathbb R\) 代数同构。设 \(M\) 为带有复共轭作用 \(\tau\) 的集合，\(F_M\) 为它所对应的 \(\operatorname{Spec}\mathbb R\) 上的平展层。写出覆盖 \(\operatorname{Spec}\mathbb C\to\operatorname{Spec}\mathbb R\) 的两条截面限制映射，由层条件计算全局截面。
\end{exercise}
\begin{solution}
所写公式双加且在 \(\mathbb R\) 上平衡，并保持乘法与单位，因而诱导代数同态。把张量积看成由第一因子给出标量作用的 \(\mathbb C\) 代数，其 \(\mathbb C\) 基为 \(1\otimes1,1\otimes i\)，两者分别映到 \((1,1),(i,-i)\)。这两个像线性无关且生成 \(\mathbb C^2\)，故映射是同构。也可用幂等元
\[
e_1=\dfrac12(1\otimes1-i\otimes i),\qquad
e_\tau=\dfrac12(1\otimes1+i\otimes i)
\]
写出逆映射 \((a,b)\mapsto(a\otimes1)e_1+(b\otimes1)e_\tau\)。

覆盖的交叠因而是两份 \(\operatorname{Spec}\mathbb C\)。两条环映射 \(w\mapsto w\otimes1\)、\(w\mapsto1\otimes w\) 在上述同构下分别为 \(w\mapsto(w,w)\)、\(w\mapsto(w,\overline w)\)。于是 \(F_M\) 的层条件是
\[
F_M(\operatorname{Spec}\mathbb R)\longrightarrow M
\rightrightarrows M\times M,
\qquad m\longmapsto(m,m),\quad m\longmapsto(m,\tau m).
\]
其等化子恰是 \(M^{\langle\tau\rangle}\)。因此这里的全局截面条件与 Galois 固定元素条件相同。
\end{solution}

\begin{exercise}\label{b4:ex:E-real-square-sections}
设 \(n\ge1\)，\(F\) 为 \(\operatorname{Spec}\mathbb R\) 上由复共轭作用的 \(\mathbb C^\times\) 所对应的交换群平展层，\(\mu_n\) 为 \(n\) 次单位根层。计算幂映射 \(F\xrightarrow{(-)^n}F\) 在全局截面上的余核。对 \(a\in\mathbb R^\times\)，选 \(z\in\mathbb C^\times\) 使 \(z^n=a\)，证明 \(\tau\mapsto\overline z/z\) 给出 \(a\) 的 Kummer 连接类，且此类为零当且仅当 \(a\) 有实 \(n\) 次根。由此计算 \(H^1_{\mathrm{\acute et}}(\operatorname{Spec}\mathbb R,\mu_n)\)。
\end{exercise}
\begin{solution}
全局截面上的映射为 \(r\mapsto r^n\)。奇数 \(n\) 时每个非零实数都有实 \(n\) 次根，余核为零群；偶数 \(n\) 时像为正实数群，符号映射给出余核 \(\mathbb R^\times/(\mathbb R^\times)^n\cong C_2\)。

令 \(G_{\mathbb R}=\{1,\tau\}\)，其中 \(\tau\) 为复共轭，定义 \(c(1)=1\)、\(c(\tau)=\overline z/z\)。因 \(a\) 为实数，\(c(\tau)^n=\overline a/a=1\)，故其值在 \(\mu_n\) 中；又有 \(c(\tau)\tau(c(\tau))=1\)，这正是 \(C_2\) 的乘法余圈条件。连接映射将一个不变截面的提升取群作用之差，在乘法记号下就是 \(g\mapsto g(z)/z\)，所以 \([c]\) 为所求连接类。若将提升换成 \(z\xi\)、\(\xi\in\mu_n\)，余圈便乘上 \(g(\xi)/\xi\)，不改变上同调类。

若 \([c]=0\)，则存在 \(\xi\in\mu_n\) 使 \(c(\tau)=\tau(\xi)/\xi\)，于是 \(z/\xi\) 被复共轭固定，且 \((z/\xi)^n=a\)，给出实根。反之，有实根便可选它作提升，使余圈恒为一。幂映射在 \(\mathbb C^\times\) 上满射，而\cref{b4:thm:C-hilbert90}给出 \(H^1(G_{\mathbb R},\mathbb C^\times)=0\)。Kummer 长正合列因此把上述余核识别为 \(H^1(G_{\mathbb R},\mu_n)\)，再由\eqref{b4:eq:C-etale-field}得到所求平展上同调：奇数 \(n\) 时为零群，偶数 \(n\) 时为 \(C_2\)。
\end{solution}
